\subsection{多项式环的同调维数}

引理 2. --- 设 \(\rho: \mathbf{A} \to \mathbf{A}'\) 为环同态，M 为 A-模， \(\mathbf{M}'\) 为 \(\mathbf{A}'\) -模。若 A-模 \(\mathbf{A}_d'\) 是平坦的，则我们有 \(\mathrm{dp}_{\mathbf{A}'}(\mathbf{A}' \otimes_{\mathbf{A}} \mathbf{M}) \leqslant \mathrm{dp}_{\mathbf{A}}(\mathbf{M})\) 。若 A-模 \(\mathbf{A}_s'\) 是投射的，则我们有 \(\mathrm{dp}_{\mathbf{A}}(\mathbf{M}') \leqslant \mathrm{dp}_{\mathbf{A}'}(\mathbf{M}')\) .

第一个断言在 \(\mathrm{dp}_{\mathbf{A}}(\mathbf{M}) = \pm \infty\) 时是显然的；若 \(\mathrm{dp}_{\mathbf{A}}(\mathbf{M}) = n \in \mathbb{N}\) ，则存在 A-模的正合序列

\[
0 \rightarrow \mathrm{P} _ {n} \rightarrow \mathrm{P} _ {n - 1} \rightarrow \dots \rightarrow \mathrm{P} _ {0} \rightarrow \mathrm{M} \rightarrow 0
\]

其中 \(P_{i}\) 是投射的；A'-模的序列

\[
0 \rightarrow \mathrm{A} ^ {\prime} \otimes_ {\mathrm{A}} \mathrm{P} _ {n} \rightarrow \mathrm{A} ^ {\prime} \otimes_ {\mathrm{A}} \mathrm{P} _ {n - 1} \rightarrow \dots \rightarrow \mathrm{A} ^ {\prime} \otimes_ {\mathrm{A}} \mathrm{P} _ {0} \rightarrow \mathrm{A} ^ {\prime} \otimes_ {\mathrm{A}} \mathrm{M} \rightarrow 0
\]

是正合的，因为 \(\mathbf{A}_d'\) 是平坦的，且 \(\mathbf{A}'\) -模 \(\mathbf{A}' \otimes_{\mathbf{A}} \mathrm{P}_i\) 是投射的（II，第 89 页，推论）；因此 \(\mathrm{dp}_{\mathbf{A}'}(\mathbf{A}' \otimes_{\mathbf{A}} \mathbf{M}) \leqslant n = \mathrm{dp}_{\mathbf{A}}(\mathbf{M})\) 。第二个断言在 \(\mathrm{dp}_{\mathbf{A}'}(\mathbf{M}') = \pm \infty\) 时是显然的；若 \(\mathrm{dp}_{\mathbf{A}'}(\mathbf{M}') = m \in \mathbf{N}\) ，则存在 \(\mathbf{A}'\) -模的一个正合序列

\[
0 \to \mathrm{P} _ {m} ^ {\prime} \to \mathrm{P} _ {m - 1} ^ {\prime} \to \dots \to \mathrm{P} _ {0} ^ {\prime} \to \mathrm{M} ^ {\prime} \to 0,
\]

其中 \(P_{i}^{\prime}\) 是投射的；底层的 A-模序列是正合的。另一方面，每个 \(P_{i}^{\prime}\) 是某个模 \(A_{s}^{\prime(I)}\) 的 A'-子模的一个直和项，因此是投射的 A-模；所以我们有 \(\mathrm{dp}_{\mathbf{A}}(\mathbf{M}^{\prime}) \leqslant m = \mathrm{dp}_{\mathbf{A}^{\prime}}(\mathbf{M}^{\prime})\) .

引理 3. --- 设 A 是交换的，并令 M 为一个 A{[}X{]}-模。

\begin{enumerate}
\def\labelenumi{\alph{enumi})}
\item
  我们有 \(\mathrm{dp}_{\mathbf{A}}(\mathbf{M}) \leqslant \mathrm{dp}_{\mathbf{A}[\mathbf{X}]}(\mathbf{M}) \leqslant \mathrm{dp}_{\mathbf{A}}(\mathbf{M}) + 1\) .
\item
  如果位似 \(X_{M}\) 是单射的，我们有 dp \(_{A}\) (M/XM) ≤ dp \(_{A[X]}\) (M)。
\item
  若 \(\mathbf{XM} = 0\) ，我们有 \(\mathrm{dp}_{\mathbf{A}}(\mathbf{M}) + 1 = \mathrm{dp}_{\mathbf{A}[\mathbf{X}]}(\mathbf{M}).\)
\item
  我们有一个 A{[}X{]}-模的正合序列（III，第 106 页及 VII，第 5 节，第 1 号）
\end{enumerate}

\[
0 \rightarrow \mathrm{A} [ \mathrm{X} ] \otimes_ {\mathrm{A}} \mathrm{M} \rightarrow \mathrm{A} [ \mathrm{X} ] \otimes_ {\mathrm{A}} \mathrm{M} \rightarrow \mathrm{M} \rightarrow 0;
\]

断言 (a) 由 X，第 135 页，推论 2 及引理 2 得出。

\begin{enumerate}
\def\labelenumi{\alph{enumi})}
\setcounter{enumi}{1}
\tightlist
\item
  若 \(\mathrm{dp}_{\mathbf{A}[\mathbf{X}]}(\mathbf{M}) = \pm \infty\) 该断言是平凡的。如果 \(\mathbf{M}\) 是投射的且非零，则 A-模 \(\mathbf{M} / \mathbf{XM}\) 等同于 \(\mathbf{A} \otimes_{\mathbf{A}[\mathbf{X}]} \mathbf{M}\) ，因此是投射的，并且我们有 \(\mathrm{dp}_{\mathbf{A}}(\mathbf{M} / \mathbf{XM}) \leqslant 0 = \mathrm{dp}_{\mathbf{A}[\mathbf{X}]}(\mathbf{M})\) 。让我们对 \(\mathrm{dp}_{\mathbf{A}[\mathbf{X}]}(\mathbf{M}) = n\) （假设其 \(>0\) ）进行归纳推理。考虑 \(\mathbf{A}[\mathbf{X}]\) -模的一个正合序列 \(0 \to \mathbf{N} \to \mathbf{L} \to \mathbf{M} \to 0\) ，其中 \(\mathbf{L}\) 是一个自由 \(\mathbf{A}[\mathbf{X}]\) -模；将该图表应用于
\end{enumerate}

\[
\begin{array}{c} 0 \longrightarrow \mathrm{N} \longrightarrow \mathrm{L} \longrightarrow \mathrm{M} \longrightarrow 0 \\ \mathrm {X_ {N}} \Big \downarrow \quad \mathrm {X_ {L}} \Big \downarrow \quad \mathrm {X_ {M}} \Big \downarrow \\ 0 \longrightarrow \mathrm{N} \longrightarrow \mathrm{L} \longrightarrow \mathrm{M} \longrightarrow 0 \end{array}
\]

将 X，第 4 页的命题 2 应用于该图表，我们看到 \(X_{N}\) 是单射的，并且我们有正合序列

\[
0 \rightarrow \mathrm{N} / \mathrm{XN} \rightarrow \mathrm{L} / \mathrm{XL} \rightarrow \mathrm{M} / \mathrm{XM} \rightarrow 0.
\]

由于 L 在 A{[}X{]} 上是自由的，L/XL 在 A 上是自由的，且我们有

\[
\mathrm{dp} _ {\mathbf {A} [ \mathbf {X} ]} (\mathbf {N}) = n - 1, \quad \mathrm{dp} _ {\mathbf {A}} (\mathbf {M} / \mathbf {X M}) \leqslant 1 + \mathrm{dp} _ {\mathbf {A}} (\mathbf {N} / \mathbf {X N});
\]

由于由归纳假设 dp \(_{A}\) (N/XN) ≤ n - 1 ，我们推出 dp \(_{A}\) (M/XM) ≤ n，这正是所要证明的。

\begin{enumerate}
\def\labelenumi{\alph{enumi})}
\setcounter{enumi}{2}
\tightlist
\item
  该断言在 \(\mathrm{dp}_{\mathbf{A}[\mathbf{X}]}(\mathbf{M}) = \pm \infty\) 时是平凡的，并且在 \(\mathrm{dp}_{\mathbf{A}[\mathbf{X}]}(\mathbf{M}) = 0\) 时也是平凡的（这是不可能的，因为 \(\mathbf{XM} = 0\) ）。因此我们可以假设 \(\mathrm{dp}_{\mathbf{A}[\mathbf{X}]}(\mathbf{M}) = n > 0\) 。如上考虑一个正合序列 \(0 \to \mathbf{N} \to \mathbf{L} \to \mathbf{M} \to 0\) ，其中 \(\mathbf{L}\) 是一个自由 \(\mathbf{A}[\mathbf{X}]\) -模，我们得到 \(\mathbf{A}\) -模的一个正合序列
\end{enumerate}

\[
0 \rightarrow \mathrm{M} \rightarrow \mathrm{N} / \mathrm{XN} \rightarrow \mathrm{L} / \mathrm{XL} \rightarrow \mathrm{M} \rightarrow 0.
\]

由 b)，我们有 \(\mathrm{dp}_{\mathbf{A}}(\mathbf{N} / \mathbf{X}\mathbf{N}) \leqslant \mathrm{dp}_{\mathbf{A}[\mathbf{X}]}(\mathbf{N}) = \mathrm{dp}_{\mathbf{A}[\mathbf{X}]}(\mathbf{M}) - 1 = n - 1\) ；由于 \(\mathrm{dp}_{\mathbf{A}}(\mathbf{L} / \mathbf{XL}) = 0\) ，从前面的正合序列出发，应用 X，第 135 页，推论 2 两次，我们推出 \(\mathrm{dp}_{\mathbf{A}}(\mathbf{M}) \leqslant n - 1\) 。但由 a)，我们有

\[
\mathrm{dp} _ {\mathbf {A}} (\mathbf {M}) \geqslant \mathrm{dp} _ {\mathbf {A} [ \mathbf {X} ]} (\mathbf {M}) - 1 = n - 1,
\]

由此得 c)。

定理 1. --- 设 A 是交换的。则

\[
\mathrm{dh} (\mathbf {A} [ \mathbf {X} ]) = \mathrm{dh} (\mathbf {A}) + 1.
\]

对每个 A{[}X{]}-模 M，我们有（引理 3）

\[
\mathrm{dp} _ {\mathbf {A} [ \mathbf {X} ]} (\mathbf {M}) \leqslant \mathrm{dp} _ {\mathbf {A}} (\mathbf {M}) + 1 \leqslant \mathrm{dh} (\mathbf {A}) + 1
\]

因此 \(\mathrm{dh}(\mathbf{A}[\mathbf{X}]) \leqslant \mathrm{dh}(\mathbf{A}) + 1\) ；反之，若 \(\mathbf{M}\) 是一个 A-模，令 \(\overline{\mathbf{M}}\) 为 \(\mathbf{A}[\mathbf{X}]\) -模，它由赋予 \(\mathbf{M}\) 使 \(\mathbf{X}\mathbf{M} = 0\) 的结构而得到，则（引理 3）

\[
\mathrm{d} p _ {\mathrm{A}} (\mathrm{M}) = \mathrm{d} p _ {\mathrm{A} [ \mathrm{X} ]} (\overline {{\mathrm{M}}}) - 1 \leqslant \mathrm{d} h (\mathrm{A} [ \mathrm{X} ]) - 1,
\]

因此 dh(A) ≤ dh(A{[}X{]}) - 1.

推论 1. --- 设 A 是交换的。我们有

\[
\mathrm{dh} (\mathrm{A} [ \mathrm{X} _ {1}, \dots , \mathrm{X} _ {n} ]) = \mathrm{dh} (\mathrm{A}) + n.
\]

这由定理对 n 进行归纳得出。

推论 2. --- 设 K 是一个交换域（相应地，一个主理想环* 或一个非域的戴德金环*）。则 dh(K{[}X₁, \ldots, Xₙ{]}) 等于 n（相应地，n + 1）。这由 dh(K) = 0（相应地，dh(K) = 1）这一事实得出。

